% Class-compatibility preamble retained for a reader excerpt of
%   Dave Witte Morris, "A Cohomological Proof that Real Representations of
%   Semisimple Lie Algebras Have Q-Forms", SIGMA 11 (2015) 034,
%   DOI 10.3842/SIGMA.2015.034, CC BY 4.0.
%
% The source pins the publisher class sigma, which is not part of the delivery
% baseline, so this file supplies the amsart baseline and then copies the author's
% own declaration block mechanically and unchanged:
%   - pinned-original-source.tex lines 11-31: the author's DeclareMathOperator
%     list, the equation numbering and the author's newtheorem declarations;
%   - pinned-original-source.tex lines 35-43, which the author writes immediately
%     after \begin{document} in the publisher layout: the author's own \Res,
%     \ResOne and \Lie macros and \allowdisplaybreaks.
% The publisher header commands of the sigma class (\FirstPageHeading,
% \ArticleName and the like) and the macro \PaperNumber they require are not
% reproduced, because this excerpt carries neither the journal front matter nor
% anything else that uses them.
% No mathematical macro, statement, symbol or assumption is introduced, changed
% or removed by this file.  The author's full DeclareMathOperator list of pinned
% source lines 11-19 is carried, including GL, SL, SO and SU: these four are not
% supplied by the delivery baseline, so leaving them out makes the author's own
% notation (for instance GL_n(Q) in the corollary on lattices) undefined.
\documentclass[10pt,a4paper]{amsart}
\usepackage{amssymb,amsmath,amsthm}
\usepackage{enumerate}
\usepackage{hyperref}

\DeclareMathOperator{\Cliff}{Clif\/f}
\DeclareMathOperator{\Gal}{Gal}
\DeclareMathOperator{\GL}{GL}
\DeclareMathOperator{\Inn}{Inn}
\DeclareMathOperator{\Mat}{Mat}
\DeclareMathOperator{\Out}{Out}
\DeclareMathOperator{\SL}{SL}
\DeclareMathOperator{\SO}{SO}
\DeclareMathOperator{\SU}{\mathbf{SU}}
\numberwithin{equation}{section}

\newtheorem{prop}{Proposition}[section]
\newtheorem{cor}[prop]{Corollary}
\newtheorem{lem}[prop]{Lemma}
{\theoremstyle{definition}
\newtheorem{defn}[prop]{Definition}
\newtheorem{notation}[prop]{Notation}
\newtheorem{rem}[prop]{Remark}
\newtheorem{ques}[prop]{Question}
\newtheorem{eg}[prop]{Example}
\newtheorem{erratum}[prop]{Erratum}}

\newcommand{\Res}[1]{\mathop{\mathbf{R}_{#1}}}
\newcommand{\ResOne}[1]{\mathop{\mathbf{R}^{(1)}_{#1}}}
\newcommand{\Lie}[1]{\mathfrak{\lowercase{#1}}}
\allowdisplaybreaks
\usepackage{cleveref}
\crefname{originallemma}{original lemma}{original lemmas}
\Crefname{originallemma}{Original lemma}{Original lemmas}
\makeatletter
\newlabel{SameFixed}{{2.3}{source}{}{sourcepointer-1}{}}
\newlabel{SameFixed@cref}{{[originallemma][0][]2.3}{[1][0][]source}{}{}{}}
\makeatother
\title{R-universality of the absolutely simple algebras of type E6: strong innerness or splitting over an imaginary quadratic extension}
\author{Dave Witte Morris}
\date{doi:10.3842/SIGMA.2015.034; CC BY 4.0}
\begin{document}
\maketitle
\noindent\textit{Editorial note.} All delivered mathematical text is the author's own, copied from the pinned publisher source sources/algebra-fl02-a-cohomological-proof-that-real-re/pinned-original-source.tex; no statement or proof step was rewritten, reordered or removed.\par\medskip
\noindent\textit{Editorial note.} Preamble provenance: the source pins the publisher class sigma, which is not part of the delivery baseline, so this card uses sources/algebra-fl02-a-cohomological-proof-that-real-re/retained-preamble.tex, which supplies the amsart baseline and then carries the author's own declarations from the pinned source (its header records the exact line ranges). The author's full DeclareMathOperator list is carried, including GL, SL, SO and SU, because the delivery baseline does not provide those four operator names. No macro was invented and no mathematical symbol was changed.\par\medskip
\noindent\textit{Editorial note.} Bibliography provenance: the citations of the counted proof are resolved with sources/algebra-fl02-a-cohomological-proof-that-real-re/bibliography.tex, an excerpt of the author's own thebibliography environment (pinned source lines 954-1016) copied character for character with all twelve entries kept; no entry was rewritten, reordered, merged, split or invented and no citation key was changed.\par\medskip
\noindent\textit{Editorial note.} External references are navigation pointers only; every deliberately kept step of the author's argument stays inside the counted spans.\par\medskip
\noindent\textit{Editorial note.} Selection provenance: this card was opened under the file-based control contract (CONTROL.md), with the mechanical pre-card gate run on its statement and proof interval before the card was written.\par\medskip
\noindent\textit{Editorial note.} The counted proof refers to the author's proposition labelled Beta=D; because that label contains an equals sign and the delivery format accepts only a restricted label alphabet, the proposition's statement is delivered verbatim as a context span instead of being declared as a navigation pointer.\par\medskip
\noindent\textit{Allowed standard prerequisites (editorial).}\begin{itemize}
\item The classification of forms of type E6 by their Tits index, and the notion of being strongly inner. Used inside the counted proof at source lines 796-800 and 818-820, attributed by the author to section 6.4.5 of the cited paper Tits-RepLinIrred. Supplies the list of strongly inner forms that the criteria are stated in terms of.
\item The fact that the centre of a simply connected almost simple group of type E6 is cyclic of prime order. Used at the start of the counted proof (source lines 793-795). Reduces the criterion to a single Tits algebra, which is what makes the two alternatives exhaustive.
\end{itemize}\medskip
\ifdefined\hypertarget\hypertarget{sourcepointer-1}{}\fi
\noindent\textit{Source reference.} Original lemma 2.3: sources/algebra-fl02-a-cohomological-proof-that-real-re/pinned-original-source.tex, source lines 471--471. Navigation pointer only: the author's lemma on the minimal inner extension; it has its own card in this lane and its content is declared as a named standard prerequisite.\par\medskip
\begin{prop}
Let $\Lie G$ be an absolutely simple Lie algebra over~$\mathbb{Q}$ of type~$E_6$.
Then $\Lie G$ is $\mathbb{R}$-universal if and only if either
\renewcommand{\labelenumi}{{\rm (\theenumi)}}
\begin{enumerate}\itemsep=0pt
\item $\Lie G$ is strongly inner $($that is, of type ${}^1 E^{28}_{6,2}$ or ${}^1 E^{0}_{6,6})$, or
\item $\Lie G$ is an outer form that splits over an imaginary quadratic extension of~$\mathbb{Q}$.
\end{enumerate}
\end{prop}

\begin{prop}[\protect{\cite[Corollary~3.5, \S~4.2, and Lemma~7.4]{Tits-RepLinIrred}}]\label{Beta=D}
Let~$L$ be the center of the algebra $D_{\Lie G, F}(\lambda):= \mathrm{End}_{\Lie G}({}^F \rho_\lambda)$ $($which,
by Schur's lemma, is a~division algebra$)$.
Then:
\renewcommand{\labelenumi}{{\rm (\theenumi)}}
\begin{enumerate}\itemsep=0pt
\item
\label{Beta=D-invt}
$\Gal(\overline{F}/L) = \{ \sigma \in \Gal(\overline{F}/F) \mid \sigma(\lambda) = \lambda \}$ $($for the $*$-action
of $\Gal(\overline{F}/F))$, and
\item $\beta_{\Lie G, L}(\lambda) = [D_{\Lie G, F}(\lambda)]$, after identifying $H^2(L; \boldsymbol\mu)$ with the
Brauer group of~$L$.
\end{enumerate}
\end{prop}

\begin{proof}
Assume, f\/irst, that $\Lie G$ is inner.
Since the center of any simply connected, almost simple group of type~$E_6$ is cyclic of prime order (namely, it is of
order~$3$), we see that $\Lie G$ is $\mathbb{R}$-universal if and only if either $\Lie G$ is strongly inner
(over~$\mathbb{Q}$) or $\Lie G \otimes_\mathbb{Q} \mathbb{R}$ is not strongly inner.
From~\cite[\S~6.4.5]{Tits-RepLinIrred}, we see that Lie algebras of type ${}^1 E^{28}_{6,2}$ or ${}^1 E^{0}_{6,6}$ are
always strongly inner, and those of type ${}^1 E^{16}_{6,2}$ are never strongly inner.
In addition, a~Lie algebra of type ${}^1 E^{78}_{6,0}$ over~$\mathbb{Q}$ cannot be strongly inner~\cite[Propositions~4
and~5]{Tits-StronglyInner}.
Since $\Lie G \otimes_\mathbb{Q} \mathbb{R}$ must be of type ${}^1 E^{28}_{6,2}$ or ${}^1 E^{0}_{6,6}$, it is strongly
inner.
Hence, $\Lie G$ is $\mathbb{R}$-universal if and only if it is strongly inner.

Suppose, now, that $\Lie G$ is outer, and let~$L$ be the unique quadratic extension of~$\mathbb{Q}$ over which $\Lie
G$ is inner.
From Lemma~\ref{SameFixed}, we know that if $\Lie G$ is $\mathbb{R}$-universal, then it must remain outer over~$\mathbb{R}$,
so~$L$~is an imaginary extension of~$\mathbb{Q}$.
Let~$\lambda$ be a~weight that represents a~nontrivial element of~$Z^*$.
Then~$\lambda$ is not f\/ixed by the $*$-action, so $D_{\Lie G, \mathbb{R}}(\lambda) = D_{\Lie G, \mathbb{C}}(\lambda) =
\mathbb{C}$ and $D_{\Lie G, \mathbb{Q}}(\lambda) = D_{\Lie G, L}(\lambda)$ (see Proposition~\ref{Beta=D}).
So $\Lie G$ is $\mathbb{R}$-universal if and only if $\Lie G \otimes_\mathbb{Q} L$ is strongly inner.

Since~$L$ is an imaginary extension, we know that $\Lie G$ splits at the inf\/inite place of~$L$.
Then, by inspection of the possible Tits indices of type ${}^1 E_6$ over a~nonarchimedean local
f\/ield~\cite[p.~58]{Tits-Classification}, we see that the central vertex of the Tits index is circled at every place, so
it must be circled over~$L$~\cite[Satz~4.3.3]{Harder-BerichtNeuere}.
Therefore, $\Lie G$ must be either split or of type ${}^1 E_{6,2}^{16}$ over~$L$.
From~\cite[\S~6.4.5]{Tits-RepLinIrred}, we see that the form of type ${}^1 E_{6,2}^{16}$ is not strongly inner.
Therefore, $\Lie G$ is $\mathbb{R}$-universal if and only if $\Lie G \otimes_\mathbb{Q} L$ is split.
\end{proof}

%% Bibliography excerpt retained for reader excerpts of
%   Dave Witte Morris, "A Cohomological Proof that Real Representations of Semisimple Lie Algebras Have Q-Forms",
%   SIGMA 11 (2015) 034, DOI 10.3842/SIGMA.2015.034.
%%
%% Every entry below is copied character for character from the author's own
% thebibliography environment of the pinned publisher source (lines 954-1016).
%% No entry was rewritten, reordered, merged, split or invented, and no citation
% key was changed.  All 12 entries are kept so that any citation of a delivered
%% excerpt resolves without guessing an entry.
\begin{thebibliography}{99}

\footnotesize\itemsep=0pt

\bibitem{DemazureGrothendieck-SGA3}
Demazure M., Grothendieck A., Sch\'emas en groups.~III. Structure des sch\'emas
  en groupes reductifs, S\'eminaire de g\'eom\'etrie alg\'ebrique du {B}ois
  {M}arie, 1962/1964, available at
  \url{http://library.msri.org/books/sga/sga/pdf/sga3-3.pdf}.

\bibitem{Harder-BerichtNeuere}
Harder G., Bericht \"uber neuere {R}esultate der {G}aloiskohomologie
  halbeinfacher {G}ruppen, \textit{Jber. Deutsch. Math.-Verein.} \textbf{70}
  (1968), 182--216, available at
  \url{http://www.digizeitschriften.de/en/dms/img/?PPN=PPN37721857X_0070&DMDID=dmdlog17}.

\bibitem{Kneser-LectGalCohoClassical}
Kneser M., Lectures on {G}alois cohomology of classical groups, \textit{Tata
  Institute of Fundamental Research Lectures on Mathematics}, Vol.~47, Tata
  Institute of Fundamental Research, Bombay, 1969, available at
  \url{http://www.math.tifr.res.in/~publ/ln/tifr47.pdf}.

\bibitem{KnusEtAl-BookOfInvols}
Knus M.A., Merkurjev A., Rost M., Tignol J.P., The book of involutions,
  \textit{American Mathematical Society Colloquium Publications}, Vol.~44,
  Amer. Math. Soc., Providence, RI, 1998.

\bibitem{Lam-IntroQuadForms}
Lam T.Y., Introduction to quadratic forms over f\/ields, \textit{Graduate Studies
  in Mathematics}, Vol.~67, Amer. Math. Soc., Providence, RI, 2005.

\bibitem{Morris-RRepsQForms}
Morris D., Real representations of semisimple {L}ie algebras have {${\mathbb
  Q}$}-forms, in Algebraic Groups and Arithmetic, Tata Institute of Fundamental
  Research, Mumbai, 2004, 469--490.

\bibitem{PlatonovRapinchukBook}
Platonov V., Rapinchuk A., Algebraic groups and number theory, \textit{Pure and
  Applied Mathematics}, Vol.~139, Academic Press, Inc., Boston, MA, 1994.

\bibitem{PrasadRapinchuk-Prescribed}
Prasad G., Rapinchuk A., On the existence of isotropic forms of semi-simple
  algebraic groups over number f\/ields with prescribed local behavior,
  \href{http://dx.doi.org/10.1016/j.aim.2006.01.001}{\textit{Adv. Math.}} \textbf{207} (2006), 646--660.

\bibitem{Raghunathan-ArithLattsSSGrps}
Raghunathan M.S., Arithmetic lattices in semisimple groups, \href{http://dx.doi.org/10.1007/BF02967980}{\textit{Proc.
  Indian Acad. Sci. Math. Sci.}} \textbf{91} (1982), 133--138.

\bibitem{Tits-Classification}
Tits J., Classif\/ication of algebraic semisimple groups, in Algebraic {G}roups
  and {D}iscontinuous {S}ubgroups ({P}roc. {S}ympos. {P}ure {M}ath., {B}oulder,
  {C}olo., 1965), Amer. Math. Soc., Providence, R.I., 1966, 33--62.

\bibitem{Tits-RepLinIrred}
Tits J., Repr\'esentations lin\'eaires irr\'eductibles d'un groupe r\'eductif
  sur un corps quelconque, \href{http://dx.doi.org/10.1515/crll.1971.247.196}{\textit{J.~Reine Angew. Math.}} \textbf{247} (1971),
  196--220.

\bibitem{Tits-StronglyInner}
Tits J., Strongly inner anisotropic forms of simple algebraic groups,
  \href{http://dx.doi.org/10.1016/0021-8693(90)90201-X}{\textit{J.~Algebra}} \textbf{131} (1990), 648--677.


\end{thebibliography}
\end{document}
